\section{Changing a message without confusing the decision problem}
\label{sim:policy-message-refinement}

A manager may receive a differently worded notice carrying the same information, or a notice that distinguishes two situations previously described alike. These changes have different consequences for a policy. The preceding requirement to fix message meaning does not prohibit either change. It requires the simulator to state which relation connects the old decision problem to the new one, and to identify the new result accordingly.

A consistent one-to-one renaming preserves the original distinction between every pair of messages. The finite comparison checks that renaming across every raw observation in every original source head, including heads not selected at the current frontier. It also preserves actor, decision dates, action meanings, source identities and second-stage memory of the manager's first action. Each total rule then has exactly one renamed counterpart. The first- and second-action refusal conditions have the same truth values after the key positions are permuted, so the full surviving rule sets correspond exactly. This is the same finite decision problem under different spellings; it is not an expansion of what the manager knows.

\subsection{A refined notice allows the old distinction to be recovered}

Now suppose a correction distinguishes a low message that was previously shared with another source history. The new message identifies which of those situations occurred. A manager can ignore that distinction and follow the old rule, or potentially act differently in the two situations. To make this precise, each refined message must map to one old message, consistently across all original raw observations. Combining this message map with the unchanged decision stage and own-action memory gives a map from new policy keys to old keys.

Write $U_o$ and $U_n$ for the old and new finite key domains, and $e:U_n\to U_o$ for this erasure map. Require $e$ to reach every old key. An old rule $\pi_o$ assigns pay or draw to each old key. Its lifted rule $L_e\pi_o$ chooses $\pi_o(e(u))$ at a new key $u$: the manager acts as if the additional distinction had not been received. New keys that erase to the same old key form a \emph{fiber}. A rule constant on each fiber makes the same choice at all of those keys.

For the same checked temporal frontier, let $C_o$ and $C_n$ be the rules satisfying the old and new known constraints. The comparison preserves each underlying action row and its refusal meaning, and checks that erasing the new forbidden assignments gives exactly the old forbidden assignments. Under those conditions,
\begin{equation}
 \{L_e\pi_o:\pi_o\in C_o\}
 =\{\pi_n\in C_n:\pi_n(u)=\pi_n(v)
       \text{ whenever }e(u)=e(v)\}.
 \label{eq:policy-refinement-lift}
\end{equation}
To see the first direction, a lifted rule makes the old choice at each corresponding raw observation. A refused first action or refused pair of successive actions would therefore also have been refused under the old rule. Conversely, a new rule constant on fibers defines a unique old rule because every old key has a preimage. The same correspondence of refused assignments then shows that this old rule satisfies the old constraints. The relation concerns exact finite rule sets under supplied message semantics; it assigns neither probabilities nor monetary values to information.

\subsection{Additional rule assignments need a reached-action comparison}

In the correction-specific low-message example, the old declaration permits four total rules and the refined declaration permits eight. Lifting preserves all four old rules. The other four also produce distinct sets of reached action rows across the complete selected family. Here the new distinction therefore permits additional feasible behavior under the stipulated information design. In another example, collapsing the first messages leaves no feasible rule, while restoring the refined messages permits four. Neither comparison establishes that those messages are actually available to a real manager or that any policy maximises value.

The source-unknown examples show why counting total rules is insufficient. With conflicting source heads, a refinement increases the unconstrained assignments from 32 to 2,048, but there is no selected history on which to demonstrate a new action. Every rule has the same empty reached-row signature. Under inconsistent clocks, the corresponding counts are 32 and 512, yet no admissible time assignment exists. An empty signature is a set used for comparison, not a realised or feasible company behavior. With unknown access, additional choices can be exhibited on the selected completed histories while the overall source family remains unresolved.

The comparison therefore retains both the full rule sets and their reached-row signatures. Pending branches, unknown source outcomes and inconsistent clocks keep their original meanings. The refined result has its own declaration and domain identity, even when its feasible rules correspond exactly to an older result.\footnote{Exact renaming is checked in \path{data/accounting-master/policy-domain-transport-20261027/}: 22 frontier comparisons and 12 refusal controls. Erasure and lifting are checked in \path{data/accounting-master/policy-observation-refinement-20261027/}: 17 comparisons and 10 controls, with whole-set equality and explicit reached-action comparisons. Split or merged labels are outside the renaming relation; a false refinement direction, altered source or dropped pending branch is refused. Both profiles remain finite conditional design evidence.}

This gives two precise ways to compare changed declarations. Broader changes to sources, decision times, actions or the observation language need additional arguments before earlier policy work can be reused.

